\chapter{Практические задания}

\section*{Задание 1}

Написать функцию, которая принимает целое число и возвращает первое четное число, не меньшее аргумента.

\begin{lstinputlisting}[
	caption={Задание 1},
	label={lst:t1},
	linerange={3-9},
	]{../src/main.lsp}
\end{lstinputlisting}

\section*{Задание 2}

Написать функцию, которая принимает число и возвращает число того же знака, но с модулем на 1 больше модуля аргумента.

\begin{lstinputlisting}[
	caption={Задание 2},
	label={lst:t2},
	linerange={30-35},
	]{../src/main.lsp}
\end{lstinputlisting}

\section*{Задание 3}

Написать функцию, которая принимает два числа и возвращает список из этих чисел, расположенный по возрастанию.

\begin{lstinputlisting}[
	caption={Задание 3},
	label={lst:t3},
	linerange={43-48},
	]{../src/main.lsp}
\end{lstinputlisting}

\section*{Задание 4}

Написать функцию, которая принимает три числа и возвращает Т только тогда, когда первое число расположено между вторым и третьим.

\begin{lstinputlisting}[
	caption={Задание 4},
	label={lst:t4},
	linerange={56-60},
	]{../src/main.lsp}
\end{lstinputlisting}

\section*{Задание 5}
Каков результат вычисления следующих выражений?

\begin{table}[ht]
	\begin{center}
		\begin{tabular}{|c|c|}
			\hline
			\textbf{Выражение} & \textbf{Результат} \\
			\hline
			(and 'fee 'fie 'foe)  & foe \\
			\hline
			(or 'fee 'fie 'foe) & fee \\
			\hline
			(or nil 'fie 'foe)  & fie \\
			\hline
			(and nil 'fie 'foe) & nil \\
			\hline
			(and (equal 'abc 'abc) 'yes) & yes \\
			\hline
			(or (equal 'abc 'abc) 'yes) & T \\
			\hline
		\end{tabular}
	\end{center}
\end{table}

\section*{Задание 6}

Написать предикат, который принимает два числа-аргумента и возвращает Т, если первое число не меньше второго.

\begin{lstinputlisting}[
	caption={Задание 6},
	label={lst:t6},
	linerange={79-81},
	]{../src/main.lsp}
\end{lstinputlisting}

\section*{Задание 7}

Какой из следующих двух вариантов предиката ошибочен и почему?

\begin{lstinputlisting}[
	caption={Задание 7},
	label={lst:t7},
	linerange={93-99},
	]{../src/main.lsp}
\end{lstinputlisting}

Ошибочный вариант будет $pred2$, т.к. сначала попытается вычислить положительное ли число путем использования $plusp$, что может вызвать ошибку, если x не является числом.

\section*{Задание 8}

Решить задачу 4, используя для ее решения конструкции: только IF, только COND, только AND/OR.

\begin{lstinputlisting}[
	caption={Задание 8},
	label={lst:t8},
	linerange={111-134},
	]{../src/main.lsp}
\end{lstinputlisting}

\section*{Задание 9}

Переписать функцию how-alike, приведенную в лекции и использующую COND, используя только конструкции IF, AND/OR.

\begin{lstinputlisting}[
	label={lst:t9-1},
	linerange={153-159},
	]{../src/main.lsp}
\end{lstinputlisting}

\begin{lstinputlisting}[
	caption={Задание 9 (Начало)},
	label={lst:t9-2},
	linerange={161-180},
	]{../src/main.lsp}
\end{lstinputlisting}

\begin{lstinputlisting}[
	caption={Задание 9 (Продолжение)},
	label={lst:t9-3},
	linerange={181-188},
	]{../src/main.lsp}
\end{lstinputlisting}